
\documentclass{article}
\usepackage{colortbl}
\usepackage{makecell}
\usepackage{multirow}
\usepackage{supertabular}

\begin{document}

\newcounter{utterance}

\centering \large Interaction Transcript for game `ifeval', experiment `ifeval', episode 53 with reflexion{-}llama{-}3.1{-}8b{-}instruct.
\vspace{24pt}

{ \footnotesize  \setcounter{utterance}{1}
\setlength{\tabcolsep}{0pt}
\begin{supertabular}{c@{$\;$}|p{.15\linewidth}@{}p{.15\linewidth}p{.15\linewidth}p{.15\linewidth}p{.15\linewidth}p{.15\linewidth}}
    \# & \multicolumn{2}{c}{Player} && \multicolumn{2}{c}{Game Master} \\
    \hline

    \theutterance \stepcounter{utterance}  
    & & & \multicolumn{4}{p{0.6\linewidth}}{
        \cellcolor[rgb]{0.9,0.9,0.9}{
            \makecell[{{p{\linewidth}}}]{
                \texttt{\tiny{[P1$\langle$GM]}}
                \texttt{You have attempted this same task before and did not succeed. These are the plans you wrote after those attempts.} \\
\texttt{Trial \#1: New plan: I failed to create a compelling and relatable journal entry about stress management. What I should have done differently is to start with a more specific and personal scenario that resonates with the target audience, in this case, college students. I should have also explored more effective stress management strategies, such as time{-}blocking, seeking support from friends or a therapist, and practicing self{-}compassion. Next time, I will start by brainstorming a list of potential stress management techniques and then craft a journal entry that incorporates one or two of these strategies in a realistic and engaging way. I will also make sure to use a more conversational tone and include specific details to make the entry feel more authentic.} \\
\\ 
\texttt{Trial \#2: New plan: I failed to create a compelling and relatable journal entry about stress management. What I should have done differently is to start with a more specific and personal scenario that resonates with the target audience, in this case, college students. I should have also explored more effective stress management strategies, such as time{-}blocking, seeking support from friends or a therapist, and practicing self{-}compassion. Next time, I will start by brainstorming a list of potential stress management techniques and then craft a journal entry that incorporates one or two of these strategies in a realistic and engaging way. I will also make sure to use a more conversational tone and include specific details to make the entry feel more authentic. } \\
\\ 
\texttt{To achieve this, I will start by writing a journal entry that begins with a personal anecdote about a time when I felt overwhelmed with stress, and then describe how I used a specific stress management technique to cope with it. For example, I could write about a time when I was struggling to meet a deadline for a project, and how I used time{-}blocking to break down the task into manageable chunks and stay focused. I will also make sure to include specific details about how the technique helped me, such as feeling more in control or having more energy.} \\
\\ 
\texttt{Trial \#3: New plan: I failed to create a compelling and relatable journal entry about stress management. What I should have done differently is to start with a more specific and personal scenario that resonates with the target audience, in this case, college students. I should have also explored more effective stress management strategies, such as time{-}blocking, seeking support from friends or a therapist, and practicing self{-}compassion. Next time, I will start by brainstorming a list of potential stress management techniques and then craft a journal entry that incorporates one or two of these strategies in a realistic and engaging way. I will also make sure to use a more conversational tone and include specific details to make the entry feel more authentic. } \\
\\ 
\texttt{To achieve this, I will start by writing a journal entry that begins with a personal anecdote about a time when I felt overwhelmed with stress, and then describe how I used a specific stress management technique to cope with it. For example, I could write about a time when I was struggling to meet a deadline for a project, and how I used time{-}blocking to break down the task into manageable chunks and stay focused. I will also make sure to include specific details about how the technique helped me, such as feeling more in control or having more energy.} \\
\\ 
\\ 
\texttt{Can you give me an example for a journal entry about stress management? Tell me how you come up with the example. Your entire response should contain less than 6 sentences.} \\
            }
        }
    }
     \\ \\

    \theutterance \stepcounter{utterance}  
    & \multicolumn{4}{p{0.6\linewidth}}{
        \cellcolor[rgb]{0.561,0.737,0.561}{
            \makecell[{{p{\linewidth}}}]{
                \texttt{\tiny{[P1$\rangle$GM]}}
                \texttt{To create a journal entry about stress management, I'll use the plan from Trial \#2. I'll start by brainstorming a list of potential stress management techniques and then craft a journal entry that incorporates one or two of these strategies in a realistic and engaging way. For example, I'll write about a time when I was struggling to meet a deadline for a project, and how I used time{-}blocking to break down the task into manageable chunks and stay focused. Here's a possible journal entry:} \\
\\ 
\texttt{"Today was a chaotic day, with multiple assignments due and a looming project deadline. I felt overwhelmed and anxious, but I remembered a trick I learned from a friend: time{-}blocking. I took a few minutes to break down the project into smaller tasks, scheduling specific times for each one. By doing so, I felt more in control and focused, and I was able to complete the project on time. It's amazing how a simple technique like time{-}blocking can make all the difference in managing stress."} \\
            }
        }
    }
    & & \\ \\

    \theutterance \stepcounter{utterance}  
    & & & \multicolumn{2}{p{0.3\linewidth}}{
        \cellcolor[rgb]{0.9,0.9,0.9}{
            \makecell[{{p{\linewidth}}}]{
                \texttt{\tiny{[GM$|$GM]}}
                {'length_constraints:number_sentences': {'relation': 'less than', 'num_sentences': 6}}
            }
        }
    }
    & & \\ \\

    \theutterance \stepcounter{utterance}  
    & & & \multicolumn{2}{p{0.3\linewidth}}{
        \cellcolor[rgb]{0.9,0.9,0.9}{
            \makecell[{{p{\linewidth}}}]{
                \texttt{\tiny{[GM$|$GM]}}
                \texttt{game\_result = LOSE} \\
            }
        }
    }
    & & \\ \\

    \theutterance \stepcounter{utterance}  
    & & \multicolumn{4}{p{0.6\linewidth}}{
        \cellcolor[rgb]{0.95,0.95,0.95}{
            \makecell[{{p{\linewidth}}}]{
                \texttt{\tiny{[P1$\leftrightarrow$P]}}
                \texttt{Post{-}episode Reflexion (not sent to GM):} \\
\texttt{New plan: I failed to create a compelling and relatable journal entry about stress management. What I should have done differently is to start with a more specific and personal scenario that resonates with the target audience, in this case, college students. I should have also explored more effective stress management strategies, such as time{-}blocking, seeking support from friends or a therapist, and practicing self{-}compassion. Next time, I will start by brainstorming a list of potential stress management techniques and then craft a journal entry that incorporates one or two of these strategies in a realistic and engaging way. I will also make sure to use a more conversational tone and include specific details to make the entry feel more authentic. } \\
\\ 
\texttt{To achieve this, I will start by writing a journal entry that begins with a personal anecdote about a time when I felt overwhelmed with stress, and then describe how I used a specific stress management technique to cope with it. For example, I could write about a time when I was struggling to meet a deadline for a project, and how I used time{-}blocking to break down the task into manageable chunks and stay focused. I will also make sure to include specific details about how the technique helped me, such as feeling more in control or having more energy.} \\
            }
        }
    }
    & \\ \\

\end{supertabular}
}

\end{document}
